\documentclass[dvipdfmx,autodetect-engine]{jsarticle}
\usepackage{amsmath,amssymb,bm}
\usepackage{enumitem}
\usepackage[dvipdfmx]{hyperref}

\title{Zhao and Shao (2015) 注釈付き日本語読解ノート}
\author{読解対象：Jiwei Zhao and Jun Shao, JASA 110(512), 1577--1590}
\date{}

\newcommand{\indep}{\mathrel{\perp\!\!\!\perp}}
\newcommand{\E}{\mathbb{E}}
\newcommand{\Prb}{\mathbb{P}}
\newcommand{\cue}[1]{\par\smallskip\noindent\textbf{原文の短い手がかり：}#1\par\smallskip}
\newcommand{\note}[1]{\par\smallskip\noindent\textbf{注釈：}#1\par\smallskip}
\newcommand{\connect}[1]{\par\smallskip\noindent\textbf{欠測IVメモへの接続：}#1\par\smallskip}

\begin{document}
\maketitle

\section*{書誌情報と読解方針}

本ノートは，Zhao and Shao (2015) ``Semiparametric Pseudo-Likelihoods in Generalized Linear Models With Nonignorable Missing Data'' を，欠測IV研究の観点から日本語で読み下すための注釈付き資料である。掲載誌は \textit{Journal of the American Statistical Association}, 110(512), 1577--1590 である。

本ノートは原論文の逐語的な全文翻訳ではなく，AbstractからAppendixとReferencesまでの論理展開を順番に追えるようにした詳細な日本語パラフレーズである。原文確認用の英語は短い専門句に限る。数式・仮定・定理・アルゴリズム・表の意味は，読解に必要な範囲で再構成して説明する。

\subsection*{用語対応表}

\begin{center}
\begin{tabular}{ll}
nonignorable missing data & 無視できない欠測，非無視可能欠測 \\
GLM & 一般化線形モデル \\
pseudo-likelihood & 疑似尤度 \\
instrumental variable & 操作変数，欠測IV，shadow variable \\
identifiability & 識別可能性 \\
response probability & 応答確率，観測確率 \\
canonical link & 正準リンク \\
asymptotic covariance & 漸近共分散 \\
concave--convex procedure & 凹凸分解型反復法 \\
\end{tabular}
\end{center}

\connect{d'Haultfoeuille (2010) は欠測確率をノンパラメトリックに識別する枠組みである。一方，Zhao and Shao (2015) は $Y\mid X$ をGLMとしてパラメトリック化し，欠測メカニズム $[R\mid Y,U]$ を指定せず，$R\indep Z\mid (Y,U)$ を使ってGLMパラメータを識別・推定する。}

\section{Abstract}

\cue{GLM / nonignorable missing data / pseudo-likelihood}

要旨は，本論文の対象を「応答変数と一部共変量に欠測があり，欠測メカニズムが非無視可能かつ未指定であるGLM」と定める。著者らは，操作変数を使う疑似尤度アプローチにより，欠測メカニズムを完全にモデル化せずに，GLMの未知パラメータを識別・推定する。

貢献は四つある。第一に，識別条件を明示する。第二に，疑似尤度最大化推定量の漸近性を示す。第三に，場合によっては漸近共分散行列とその推定量を導く。第四に，非凹な疑似尤度最大化を二つの凹最大化問題へ分解する二段階反復アルゴリズムを提案する。数値例としてシミュレーションと綿工場労働者データの応用が示される。

\note{この論文の「IV」は，処置効果の通常IVではなく，$Y$ とは関連するが，$Y$ と $U$ を条件づけると欠測指標 $R$ に追加情報を持たない変数である。}

\section{1. Introduction}

導入部は，標本調査や生物医学研究で欠測が一般的であり，MARでない欠測にMAR用の方法を適用すると推定バイアスと誤った推論が生じる，という問題意識から始まる。非無視可能欠測では，欠測メカニズムが欠測値そのものに依存しうるため，観測データだけで母集団パラメータを識別することが難しい。

既存研究として，Greenlees et al. の完全パラメトリック最尤，Ibrahim et al. のGLMMとMonte Carlo EM，Qin et al. の経験尤度，Kott and Chang の較正重み，Kim and Yu の指数傾斜モデル，Tang et al. の疑似尤度が整理される。完全パラメトリック法はモデル指定に敏感であり，完全ノンパラメトリック法は非無視可能欠測では一般に識別不能であるため，半パラメトリックな妥協が重要になる。

本論文は Tang, Little and Raghunathan (2003) の疑似尤度を引き継ぐが，重要な違いがある。Tangらは，$Y$ を条件づけると欠測指標 $R$ が全共変量 $X$ と独立であると仮定した。Zhao and Shao は $X$ を $U$ と $Z$ に分け，$R\indep Z\mid (Y,U)$ のみを仮定する。これにより，$U$ の効果を推定・検定しつつ，$Z$ を識別用の操作変数として使える。

この分解は実務的にも重要である。$Z$ は $Y$ と関連している必要があるため，$Z$ 自身の効果を検定するには適さない。一方，$U$ に入れた共変量の効果は識別できれば検定可能である。したがって，どの変数をIVにして，どの変数を通常共変量として残すかが設計上の核心になる。

導入部はさらに，GLMごとに識別条件が異なることを強調する。正規回帰で知られていた「IVのサポート点が三つ必要」という結論は一部GLMには残るが，ロジスティックやポアソンなどでは分散パラメータの有無により必要条件が変わる。

\section{2. Methodology}

\subsection{2.1 Models}

関心は，応答変数 $Y$ が共変量 $X=(U,Z)$ の下でGLMに従うときのパラメータ $\theta$ の推定である。密度は指数型分布族として
\[
p(y\mid x;\theta)=\exp\left\{\frac{y\eta-b(\eta)}{\phi}+c(y;\phi)\right\}
\]
と書かれ，リンク関数は
\[
g\{\mu(\eta)\}=\alpha+\beta_u^\top u+\beta_z^\top z
\]
である。$\phi$ は既知または未知の分散パラメータで，$\theta=(\alpha,\beta_u,\beta_z,\phi)$ が未知である。

扱うGLMには，正規線形回帰，ポアソン回帰，二項ロジスティック回帰，ガンマ回帰，逆ガウス回帰が含まれる。正準リンクの場合は $g=\mu^{-1}$ と書かれ，識別条件の記述が簡単になる。

データは $N$ 人の独立同分布標本で，$R=1$ のときだけ $Y$ が観測される。$X$ は基本設定では完全観測である。後の第5節では，$U$ の一部にも欠測がある場合へ拡張する。

\subsection{2.2 Identifiability}

\cue{$[R\mid Y,U,Z]=[R\mid Y,U]$}

非無視可能欠測では，$[R\mid Y,X]$ と $[Y\mid X]$ をともにパラメトリックに置いても識別できない場合がある。そこで本論文は，欠測メカニズムに関して
\[
[R\mid Y,U,Z]=[R\mid Y,U]
\]
を仮定する。これは，$Y$ と $U$ が与えられれば，$Z$ は欠測確率に追加的な情報を持たないという条件である。ただし，$R$ と $Z$ が無条件に独立であるとは言っていない。

この条件の下では，完全ケースにおける $[Z\mid Y,U,R=1]$ が母集団の $[Z\mid Y,U]$ と一致する。Bayes公式を使うと，
\[
[Z\mid Y,U]
=
\frac{p(Y\mid U,Z;\theta)[Z\mid U]}
{\int p(Y\mid U,z;\theta)[z\mid U]\,dz}
\]
と表される。したがって，完全ケースから観測される $[Z\mid Y,U]$ を用いて，GLMパラメータ $\theta$ を識別する。

\begin{description}[leftmargin=3.2em]
\item[Lemma 1] 密度のうち $Z$ を含む部分が $\delta(\theta)$ によって変化し，異なる $\delta$ が $(Y,U)$ だけの関数差に吸収されないなら，$\delta(\theta)$ は $[Z\mid Y,U]$ から識別される。
\end{description}

補題1はGLMに限らない一般結果である。要点は，$[Z\mid Y,U]$ の中でパラメータが $Z$ 方向の形を変えるなら，その部分は識別できるということである。

GLMに対しては，サポート点 $u_0,\ldots,u_q$ と $z_0,\ldots,z_k$ を取り，それらで作られる行列 $M$ のフルランク性を識別条件とする。この行列は，リンク関数や cumulant 関数 $b$ の差分を並べたもので，GLMパラメータの局所的な変化が $[Z\mid Y,U]$ に十分異なる影響を与えるかを測る。

\begin{description}[leftmargin=3.2em]
\item[Theorem 1] $2k(q+1)\ge p+2$ かつ $M$ がフルランクなら，$\theta=(\alpha,\beta_u,\beta_z,\phi)$ は識別される。$\phi$ が既知なら，最後の列を除いた $M_0$ のフルランク性と $2k(q+1)\ge p+1$ でよい。
\end{description}

重要な含意は，$\beta_{z0}\neq0$ が必要である点である。$Z$ が $Y$ に効かなければ，$[Z\mid Y,U]$ はパラメータに依存せず，識別情報を持たない。つまりIVは「欠測には直接効かない」が「アウトカムには効く」必要がある。

\begin{description}[leftmargin=3.2em]
\item[Corollary 1] $X=Z$ で正準リンクを使う場合，$\phi$ が未知なら，$Z$ のサポートに少なくとも三つの点があり $\beta_z\neq0$ なら全パラメータが識別される。$\phi$ が既知なら，$Z$ が非定数で $\beta_z\neq0$ ならよい。
\item[Corollary 2] $U$ も存在する場合，$Z$ が二値でも，$U$ が十分に豊かなサポートを持ち $\beta_u\neq0,\beta_z\neq0$ なら，未知分散パラメータを含む全パラメータが識別されうる。$\phi$ が既知なら条件はさらに弱くなる。
\end{description}

正規線形回帰やポアソン回帰では，未知分散パラメータがある場合，二値 $Z$ だけでは足りないことがある。一方，二項，ガンマ，逆ガウスなどでは，$U$ の存在が二値IVの不足を補う場合がある。したがって，識別可能性はGLMの種類，リンク，分散パラメータの有無，$U,Z$ のサポートに依存する。

\note{d'Haultfoeuille (2010) の完備性が無限次元のランク条件だとすれば，Zhao and Shao (2015) の行列 $M$ のフルランク条件は，GLMという有限次元モデルにおける具体的なランク条件である。}

\subsection{2.3 Likelihoods for Estimation}

推定では，識別式に現れる $[Z\mid Y,U]$ を使って疑似尤度を作る。問題は，$[U,Z]$ または $[Z\mid U]$ をどう扱うかである。

完全パラメトリック法では，$[U\mid Z]$ と $[Z]$ の両方にパラメトリックモデルを置く。しかし未知パラメータが増え，積分が閉形式を持たないことが多く，実装が難しい。

第一の半パラメトリック疑似尤度は，$[U\mid Z]$ だけをパラメトリックに置き，$[Z]$ は経験分布で置き換える。これは論文の式(6)であり，
\[
L(\theta)=\prod_{i=1}^n
\frac{p(y_i\mid u_i,z_i;\theta)p(u_i\mid z_i;\widehat\gamma)}
{\sum_{j=1}^N p(y_i\mid u_i,z_j;\theta)p(u_i\mid z_j;\widehat\gamma)}
\]
という形である。

第二の疑似尤度は，$[U,Z]$ 全体をカーネル密度推定で置き換える。これは式(7)である。$[U\mid Z]$ のモデル誤指定に強いが，次元が高いとカーネル推定が不安定になり効率が落ちる。

第三に，$U$ が有限カテゴリなら，$U$ で条件づけて，各カテゴリ内の $Z$ の経験分布を使う。これは式(8)であり，離散 $U$ に対する頑健な実装である。

IV選択について，論文は三つの実務指針を示す。第一に，$Z$ は $Y$ と関連していなければならない。第二に，$Z$ の次元は小さい方が，$R\indep Z\mid (Y,U)$ の破れに対して頑健である。第三に，複数のIV候補で感度分析を行うのが望ましい。

\section{3. Asymptotic Theory}

第3節は，式(6)と式(7)の疑似尤度最大化推定量の一致性と漸近正規性を扱う。式(8)は式(6)の特殊形として扱える。

\begin{description}[leftmargin=3.2em]
\item[Theorem 2] 識別条件に加えて，妨害パラメータ推定量や経験分布・カーネル推定量が適切に収束し，経験過程の一様収束条件が成り立てば，疑似尤度最大化推定量 $\widehat\theta$ は真値 $\theta_0$ に一致する。
\end{description}

式(6)では，$[U\mid Z]$ のパラメータ $\gamma$ の推定と $Z$ の経験分布を扱う。式(7)では，カーネル $K$ の次数，滑らかさ，バンド幅 $h$ に条件が必要である。$p$ 次元の $X$ でカーネルを使うため，次元が大きい場合は高次カーネルや慎重なバンド幅選択が必要になる。

\begin{description}[leftmargin=3.2em]
\item[Theorem 3] $\widehat\theta$ が一致し，正則性条件が成り立てば，$\sqrt N(\widehat\theta-\theta_0)$ は漸近正規分布に従う。式(6)では，$\gamma$ 推定と経験分布推定の影響を含む influence function 表現が得られる。式(7)にも対応するが，条件は複雑になる。
\item[Theorem 4] 式(6)の推定量について，漸近共分散行列の一貫推定量を構成できる。
\end{description}

式(7)のカーネル版では漸近共分散の明示式が複雑なため，論文は分散推定にbootstrapを推奨する。

\section{4. An Algorithm}

疑似尤度は，$p(y\mid u,z;\theta)$ の対数が凹であっても，全体としては非凹になりうる。Tang et al. はNewton型アルゴリズムを用いたが，パラメータ次元が大きいとHessianの逆行列が不安定になる。

本論文は，正準リンクGLMについて，目的関数を二つの関数の差
\[
l(\lambda;\omega)=A(\lambda;\omega)-B(\lambda;\omega)
\]
として書き，反復時点 $(\lambda_t,\omega_t)$ から二段階で更新する。

第一段階では，$\omega_t$ を固定し，$B$ を一次近似して，凹関数
\[
Q(\lambda;\omega_t)=A(\lambda;\omega_t)-\nabla_\lambda B(\lambda_t;\omega_t)^\top\lambda
\]
を最大化する。第二段階では，$\lambda_{t+1}$ を固定し，$\omega$ について凹な問題を解く。これにより各反復で目的関数が増加する。

この方法は concave--convex procedure の発想に基づく。厳密な凹性が足りない部分については，補助関数 $E$ を加えて差分を凹関数同士の差として扱う。シミュレーションでは条件がよく満たされ，安定に動作したと報告されている。

\note{実装面では，非無視可能欠測GLMの推定が「理論的には識別できるが数値的に不安定」という問題を持つため，このアルゴリズムは重要である。}

\section{5. Extension to the Case of Missing Covariates}

第5節は，$Y$ だけでなく，共変量 $U$ の一部も欠測する場合への拡張である。まず $U$ の一成分 $U_1$ が欠測する場合を考え，$R_Y$ と $R_{U_1}$ をそれぞれ $Y$ と $U_1$ の観測指標とする。

第一段階では，$U_1$ の欠測メカニズムについて
\[
[R_{U_1}\mid U,Z]=[R_{U_1}\mid U]
\]
を仮定する。つまり，$Z$ は $U_1$ に対する操作変数として働く。$[U_1\mid U_{-1},Z]$ の識別条件は第2節のGLM識別条件と同様に導ける。

第二段階では，
\[
[R_Y,R_{U_1}\mid Y,U,Z]=[R_Y,R_{U_1}\mid Y,U]
\]
を仮定し，$Z$ を $Y$ に対する操作変数として使う。複数の $U$ 成分が欠測する場合は，同じ考え方で多段階手続きを構成する。

もし $U$ の欠測がMARなら，$[U\mid Z]$ または $[U,Z]$ の推定は通常のMAR手法で行い，その後に $Y$ の非無視可能欠測に対する疑似尤度を適用すればよい。

\section{6. Simulation Studies}

シミュレーションは，提案推定量の有限標本性能を調べる。$Y,U,Z$ は一変量で，$Y$ の型は正規，二値，ポアソンの三種類，$U$ は連続または二値，$U$ の欠測あり・なしを組み合わせた12ケースである。

ケースAは正規 $Y$，ケースBは二値 $Y$，ケースCはポアソン $Y$ である。各ケースで $Y$ の欠測確率は $Y$ と $U$ に依存し，非無視可能になる。$U$ に欠測があるケースでは，$U$ の観測確率も $U$ に依存する。総標本サイズは $N=500$，完全データ率は $Y$ だけ欠測なら約80\%，$Y$ と $U$ の両方欠測なら約70\%である。

比較対象は，欠測なしの標準推定，MAR仮定の推定，提案疑似尤度(6)，必要に応じてカーネル版(7)，Tang et al. の方法，完全パラメトリック尤度である。完全パラメトリック尤度は正しく指定されれば効率的だが，欠測モデルが誤指定されるとバイアスが生じる。

表1は正規 $Y$，表2は二値 $Y$，表3はポアソン $Y$ の結果である。全体として，提案推定量のbiasは小さく，標準誤差推定もよく，95\%信頼区間の被覆率も概ね名目水準に近い。一方，MAR仮定の方法は，非無視可能欠測の下で一部パラメータに大きなバイアスと低い被覆率を示す。

式(6)と式(7)の比較では，$[U\mid Z]$ が正しく指定される場合は式(6)が効率的であり，誤指定される場合はカーネル版の式(7)がより頑健である。ただし，式(6)も完全に破綻するわけではない。

\begin{table}[tbp]
\centering
\caption{シミュレーション表1--3の読解用要約}
\begin{tabular}{p{.16\linewidth}p{.29\linewidth}p{.45\linewidth}}
\hline
表 & 対象モデル & 読解上のポイント \\
\hline
Table 1 & 正規応答 $Y$ & 欠測モデルが正しく指定された完全パラメトリック法は効率的だが，誤指定されるとバイアスが出る。提案疑似尤度は安定している。\\
Table 2 & 二値応答 $Y$ & MAR法は一部係数で被覆率が悪化する。式(6)と式(7)の効率・頑健性の違いが比較される。\\
Table 3 & ポアソン応答 $Y$ & カウントデータでも提案法のbiasと標準誤差推定は概ね良好で，GLM一般への適用可能性を示す。\\
\hline
\end{tabular}
\end{table}

\note{原論文の表は多数の数値を含むため，本ノートでは再掲せず，設計・比較対象・結論を要約表として整理した。数値確認が必要な場合は原PDFのTable 1--3を参照する。}

\section{7. Real Data Analysis}

実データ応用は，綿工場労働者の呼吸困難（dyspnoea）に関する環境研究データである。912人が1981年，1986年，1992年に調査され，それぞれの呼吸困難状態を $Y_0,Y_1,Y_2$ とする。$Y_0$ は欠測なし，$Y_1$ と $Y_2$ には欠測がある。その他の共変量は，綿粉塵曝露，性別，年齢である。

先行研究のIbrahim et al. は，欠測指標と応答ベクトルに完全パラメトリックな混合モデルを置いた。MAR仮定と非無視可能欠測仮定で結果がかなり異なり，特に時間効果の解釈に違いが出た。

本論文では，$Y=Y_2$ をアウトカムとし，$X=(\text{expd},\text{sex},\text{age},Y_0,Y_1)$ を共変量とするロジスティック回帰を考える。$Y_1$ にも欠測があるため，第5節の拡張を用いる。

IV選択では，$Z$ はアウトカムに関連し，かつ欠測メカニズムから条件付きで除外される必要がある。expdとsexは有用なIVではないと判断され，主分析では $Z=Y_0$ が選ばれる。これは，$Y_1,Y_2$ と他の共変量を条件づけると，1981年の呼吸困難状態 $Y_0$ は1992年応答の欠測に直接関係しない，という仮定である。比較のため，$Z=\text{age}$ も試す。

表4は，MAR法，提案法(6)で $Z=Y_0$，提案法(7)で $Z=Y_0$，提案法(6)で $Z=\text{age}$ の結果を比較する。共通結論として，曝露と性別の効果は有意でない。大きな違いは $Y_1$ の係数であり，MARでは小さすぎる。直観的には，1986年の呼吸困難状態は1981年の状態より1992年の状態を強く予測するはずであり，提案法の結果の方が自然である。

$Z=Y_0$ と $Z=\text{age}$ を比べると，$Y_0$ をIVにした方が合理的で標準誤差も小さい。ageはアウトカムへの効果が弱く，IVとしての関連性が不足しているため，推定が不安定になる。最後に，式(6)と式(7)の結果差はこのデータでは小さい。

\begin{table}[tbp]
\centering
\caption{実データ分析Table 4の読解用要約}
\begin{tabular}{p{.24\linewidth}p{.32\linewidth}p{.34\linewidth}}
\hline
比較 & 何を見ているか & 主要な読み取り \\
\hline
MAR vs 提案法 & 1992年呼吸困難のロジスティック回帰係数 & MARでは1986年呼吸困難 $Y_1$ の効果が過小になり，時系列的な直観と合いにくい。\\
$Z=Y_0$ & 1981年呼吸困難をIVにする設定 & 関連性があり，条件付き除外制約も比較的説明しやすいため，主分析として自然。\\
$Z=\text{age}$ & 年齢をIVにする感度分析 & 年齢のアウトカム関連性が弱く，標準誤差が大きくなり，IVとして信頼しにくい。\\
式(6) vs 式(7) & $[U\mid Z]$ モデル化版とカーネル版 & この応用では結論差は小さく，主要な違いはIV選択に由来する。\\
\hline
\end{tabular}
\end{table}

\note{著者らは，欠測仮定は検証不能であり，真の欠測メカニズムは分からないと明示する。提案法は真実を保証するものではなく，非無視可能欠測が疑われるときの有用な感度分析・代替推定として位置づけられる。}

\section{Appendix: Proofs}

付録は，本文の補題・定理・系の証明をまとめる。補題1の証明では，二つのパラメータが同じ $[Z\mid Y,U]$ を生むと仮定し，指数型の比を整理して，$Z$ を含む部分の差が $(Y,U)$ だけの関数でなければならないことを示す。補題の仮定に反するため，識別が従う。

Theorem 1の証明では，サポート点で評価した複数の識別可能な関数を作り，それらのJacobianがフルランクであれば逆関数定理によって $\theta$ が局所的に回復できることを示す。Corollary 1と2は，正準リンク・一次元 $U,Z$ のもとでこのランク条件を具体化する。

Theorem 2の証明はM推定量の一致性である。期待疑似尤度が真値で一意に最大化されることと，標本疑似尤度が一様に期待値へ近づくことを使う。Theorem 3はTaylor展開によりscore方程式を線形化し，妨害推定量と経験分布・カーネル推定の影響をinfluence functionへ吸収する。Theorem 4は，式(6)の場合に得られるinfluence function表現を標本共分散で推定できることを示す。

\section{Referencesの位置づけ}

参考文献は四つに分けられる。第一に，Little and Rubin，Robins and Ritov など欠測データと識別不能性の基礎文献。第二に，Tang et al.，Qin et al.，Kim and Yu など非無視可能欠測の半パラメトリック推定。第三に，McCullagh and Nelder などGLMの基礎。第四に，Newey，Newey and McFadden，Gong and Samaniego など疑似尤度・M推定・半パラメトリック漸近理論である。

\connect{本プロジェクトでは，Zhao and Shao (2015) は「欠測IVを使うが，欠測確率そのものをノンパラメトリックに復元するのではなく，GLMパラメータを識別・推定する」文献として位置づく。d'Haultfoeuille 型の完全ノンパラ識別と，GLMに限定した半パラメトリック疑似尤度の中間関係を整理するうえで重要である。}

\section*{全体まとめ}

Zhao and Shao (2015) の核心は，$Y$ が非無視可能に欠測しても，$Y\mid U,Z$ にGLMを仮定し，$R\indep Z\mid(Y,U)$ を満たすIV $Z$ があれば，完全ケースの $[Z\mid Y,U]$ からGLMパラメータを識別できるという点である。識別はGLMの有限次元ランク条件で特徴づけられ，推定は半パラメトリック疑似尤度で行う。$[U\mid Z]$ をモデル化する式(6)は効率的，カーネル版の式(7)は頑健であり，IV選択の感度分析が重要である。

\end{document}
